\chapter{The three invariants of a chiral algebra}
\label{ch:three-invariants}

\begin{abstract}
Three distinct numerical invariants of a chiral algebra~$\cA$
(the \emph{generator OPE pole order}~$p_{\max}(\cA)$, the
\emph{collision depth}~$k_{\max}(\cA)$, and the
\emph{shadow depth}~$r_{\max}(\cA)$) measure three different
structural features and need not coincide. This chapter formalizes
the distinction, states the relation
$k_{\max} = p_{\max} - 1$ on the fixed minimal-generator lane of the
standard landscape
(a consequence of the $\mathrm{d}\log$ absorption of the bar
propagator), and proves that $r_{\max}$ is independent
of $p_{\max}$. The $\beta\gamma$ system is the archetypal witness:
$p_{\max} = 1$, $k_{\max} = 0$, $r_{\max} = 4$.
The same separation fences these three finite pole/shadow invariants
from the scalar characteristic~$\kappa$, from the derived chiral
centre $\ChirHoch^\bullet(\cA,\cA)$, and from the bar--cobar
Koszul package.
This distinction is load-bearing for the operator-order trichotomy
of Theorem~\ref{thm:shadow-depth-operator-order}, for the
classification of the standard landscape into the four
classes~$G,L,C,M$, and for the universal-conductor
identification
(Chapter~\ref{chap:universal-conductor},
$K(\cA) = -c_{\mathrm{ghost}}(\mathrm{BRST}(\cA))$): the
shadow-depth invariant $r_{\max}$ stratifies the conductor's
domain, while $p_{\max}$ and $k_{\max}$ control the OPE-side
input data. The shadow quadrichotomy
(Chapter~\ref{chap:shadow-quadrichotomy-platonic}) reads the
$r_{\max} \in \{2, 3, 4, \infty\}$ values as the Riccati strata
$G/L/C/M$.
\end{abstract}

\section{Motivation: the \texorpdfstring{$\beta\gamma$}{beta-gamma} subtlety}
\label{sec:three-invariants-motivation}

The $\beta\gamma$ system, the chiral algebra generated by a pair
of fields $\beta(z),\gamma(z)$ of conformal weights~$1$ and~$0$
respectively, with OPE
\begin{equation}\label{eq:beta-gamma-ope}
\beta(z)\gamma(w) \;\sim\; \frac{1}{z-w},
\qquad
\gamma(z)\beta(w) \;\sim\; \frac{-1}{z-w},
\qquad
\beta\beta \;\sim\; 0,
\qquad
\gamma\gamma \;\sim\; 0
\end{equation}
illustrates the distinction between three invariants that are
often conflated in the literature. The maximal OPE pole
in~\eqref{eq:beta-gamma-ope} is a \emph{simple} pole: the
\emph{generator OPE pole order} of $\beta\gamma$ is~$1$. By the
$\mathrm{d}\log$ absorption of the bar
propagator~(Remark~\textup{\ref{rem:bar-pole-absorption}}), this
simple pole is absorbed by the $\mathrm{d}\log(z-w)$ kernel and
produces no collision residue: the \emph{collision depth} of
$\beta\gamma$ is~$0$.

Yet the $\beta\gamma$ system exhibits genuine higher-degree shadow
obstructions. The composite current $J={:}\beta\gamma{:}$ has
\[
J(z)J(w)\;\sim\;\frac{\varepsilon_{\beta\gamma}}{(z-w)^2}
       + \frac{\textup{current}}{z-w},
       \qquad \varepsilon_{\beta\gamma}\in\{\pm 1\},
\]
by the two Wick contractions of the mixed simple-pole OPE. This is
not a new generator-pair pole for the minimal generating set
$\{\beta,\gamma\}$; it is a composite channel in the ordered bar
transfer. On the standard conformal-weight family the corresponding
shadow datum is
\[
 S_3(\beta\gamma)=0,\qquad
 S_4(\beta\gamma)=-\frac{5}{12},\qquad
 S_r(\beta\gamma)=0\quad(r\geq 5),
\]
while the rank-one weight-changing scalar projection
$\mu_{\beta\gamma}$ vanishes. The \emph{shadow depth} of
$\beta\gamma$ is therefore~$4$, which places it in class~$C$ of the
shadow-depth classification.

The three invariants for $\beta\gamma$ are thus
\[
 p_{\max} = 1,
 \qquad
 k_{\max} = 0,
 \qquad
 r_{\max} = 4,
\]
and they disagree maximally. This chapter formalizes the three
invariants and proves the relations among them.

\section{Definitions}
\label{sec:three-invariants-definitions}

\begin{definition}[Generator OPE pole order]
\label{def:p-max}
\index{generator OPE pole order $p_{\max}$|textbf}
\index{p_max@$p_{\max}$|see {generator OPE pole order}}
Let $\cA$ be a chiral algebra with a chosen set of generating
fields $\{a_\alpha\}$. The \emph{generator OPE pole order}
of~$\cA$ is
\begin{equation}\label{eq:p-max-def}
p_{\max}(\cA)
\;:=\;
\max_{\alpha,\beta}\,
\Big\{\,n \geq 0
\;\Big|\;
a_\alpha(z)a_\beta(w)
\;\text{has a nonzero}\;\frac{1}{(z-w)^n}\;\text{term}\,\Big\}.
\end{equation}
Equivalently, $p_{\max}(\cA)$ is the maximal~$n$ such that some
OPE coefficient $a_{\alpha,(n-1)}a_\beta$ among
generating-field pairs is nonzero.
\end{definition}

\begin{remark}[Independence from the choice of generators]
\label{rem:p-max-generator-choice}
The invariant $p_{\max}(\cA)$ depends on the choice of generators.
For a strongly generated chiral algebra, the standard choice is the
minimal set of strong generators, and $p_{\max}$ is the maximal OPE
pole order among that minimal set. For the standard landscape
\textup{(}Heisenberg, affine Kac--Moody, Virasoro, $\Walg_N$,
lattice VOAs, $\beta\gamma$\textup{)}, the minimal generator set is
canonical up to automorphism, so $p_{\max}$ is well-defined.
\end{remark}

\begin{definition}[Collision depth]
\label{def:k-max}
\index{collision depth $k_{\max}$|textbf}
\index{k_max@$k_{\max}$|see {collision depth}}
The \emph{collision depth} of~$\cA$ is
\begin{equation}\label{eq:k-max-def}
k_{\max}(\cA)
\;:=\;
\max\Big\{\,k \geq 0
\;\Big|\;
\operatorname{Res}^{\mathrm{coll}}_{0,2}(\Theta_\cA)
\;\text{has a nonzero}\;\frac{1}{z^k}\;\text{term}\,\Big\}.
\end{equation}
The value~$k=0$ records a nonzero regular scalar term after the
$d\log$ kernel has absorbed the generator simple pole; positive
values record genuine binary collision poles.
By the $\mathrm{d}\log$ absorption of the bar propagator
\textup{(}Remark~\textup{\ref{rem:bar-pole-absorption}}\textup{)},
the collision residue of a term $1/(z-w)^n$ in the OPE contributes
to $1/z^{n-1}$ in the collision residue: the $\mathrm{d}\log(z-w)$
kernel absorbs one power. Consequently
\begin{equation}\label{eq:k-max-from-p-max}
k_{\max}(\cA)
\;=\;
p_{\max}(\cA) - 1.
\end{equation}
\end{definition}

\begin{remark}[Binary convention and regular terms]
\label{rem:kmax-regular-term}
The definition is binary and generator-relative. It reads the
two-point collision residue of the ordered bar differential on the
chosen minimal strong generators. For $\beta\gamma$ the mixed
generator OPE has a simple pole, and after $d\log$ absorption its
binary collision residue is regular; this is
$k_{\max}(\beta\gamma)=0$, not absence of the $\beta\gamma$ OPE. The
quartic class that gives $r_{\max}=4$ belongs to the all-degree
shadow tower, after normal-ordered composites enter.
\end{remark}

\begin{remark}[Why $k_{\max}$ is not redundant]
\label{rem:k-max-not-redundant}
Although $k_{\max}$ is determined by $p_{\max}$
via~\eqref{eq:k-max-from-p-max}, the two invariants measure different
things: $p_{\max}$ is a property of the chiral algebra's OPE data,
while $k_{\max}$ is a property of the bar complex at degree~$2$. The
distinction matters because $k_{\max}$ is the invariant that appears
in the GZ\textup{26} commuting-Hamiltonian operator-order formula
\textup{(}Theorem~\textup{\ref{thm:shadow-depth-operator-order}}\textup{)}:
the order of the differential operator~$H_i$ is $k_{\max} - 1$
\textup{(}for $k_{\max}\geq 1$\textup{)} or $0$
\textup{(}for $k_{\max} = 0$\textup{)}.
\end{remark}

\begin{definition}[Shadow depth]
\label{def:r-max}
\index{shadow depth $r_{\max}$|textbf}
\index{r_max@$r_{\max}$|see {shadow depth}}
The \emph{shadow depth} of $\cA$ is
\begin{equation}\label{eq:r-max-def}
r_{\max}(\cA)
\;:=\;
\max\Big\{\,r \geq 2
\;\Big|\;
\text{the shadow obstruction tower }\Theta_\cA^{\leq r}
\text{ has a nonzero~$r$-degree projection}\,\Big\}
\;\in\;\{2, 3, 4, \infty\}.
\end{equation}
Equivalently, $r_{\max}(\cA)$ is the degree at which the finite-order
engine $\Theta_\cA^{\leq r}$ first terminates: for $r_{\max}(\cA)<\infty$,
the shadow obstruction tower is nontrivial only through
degree~$r_{\max}$, and all higher-degree projections vanish.
For $r_{\max}(\cA) = \infty$, the tower is nontrivial at all degrees.
\end{definition}

\begin{remark}[Shadow depth depends on all degrees]
\label{rem:r-max-all-degrees}
Unlike $p_{\max}$ and $k_{\max}$, which are degree-$2$ invariants,
the shadow depth $r_{\max}$ involves composites and normal-ordered
products that enter at higher degrees. For $\beta\gamma$, the
generator self-OPE is simple pole only
\textup{(}$p_{\max}=1$\textup{)}, but the composite $:\!\beta\gamma\!:$
at degree~$4$ carries a nontrivial quartic contact class,
giving $r_{\max}(\beta\gamma) = 4$. In general,
$r_{\max}(\cA)$ is \emph{not} determined by the OPE pole structure
of the generators alone.
\end{remark}

\section{Relations among the three invariants}
\label{sec:three-invariants-relations}

\begin{proposition}[Relations and independence]
\label{prop:three-invariants-relations}
\ClaimStatusProvedHere
For every chiral algebra $\cA$ in the standard landscape:
\begin{enumerate}[label=\textup{(\roman*)}]
\item $k_{\max}(\cA) = p_{\max}(\cA) - 1$.
\item $r_{\max}(\cA)$ is independent of $p_{\max}(\cA)$: the pair
 $(p_{\max}, r_{\max})$ achieves values
 $(1, 4)$ \textup{(}$\beta\gamma$\textup{)},
 $(2, 2)$ \textup{(}Heisenberg\textup{)},
 $(2, 3)$ \textup{(}affine $\mathfrak{sl}_N$\textup{)},
 $(4, \infty)$ \textup{(}Virasoro\textup{)},
 $(2N, \infty)$ \textup{(}$\Walg_N$\textup{)},
 among others. In particular, $p_{\max} = 2$ does not determine
 $r_{\max}\in\{2, 3\}$, and $r_{\max} = \infty$ does not determine
 $p_{\max}$.
\item The pair $(r_{\max}, \mathrm{class})$ classifies the standard
 landscape into four shadow-depth classes:
 \begin{center}
 \renewcommand{\arraystretch}{1.2}
 \begin{tabular}{cccl}
 \textbf{Class} & $r_{\max}$ & \textbf{Examples} & \textbf{OPE pole structure} \\
 \hline
 $G$ & $2$ & Heisenberg, lattice VOAs & quadratic ($p_{\max}=2$) \\
 $L$ & $3$ & Affine $\widehat{\fg}_k$ & quadratic ($p_{\max}=2$) \\
 $C$ & $4$ & $\beta\gamma$, matrix fact. & simple ($p_{\max}=1$) \\
 $M$ & $\infty$ & Virasoro, $\Walg_N$ & quartic+ ($p_{\max}\geq 4$)
 \end{tabular}
 \end{center}
\end{enumerate}
\end{proposition}

\begin{proof}
Part~(i) is the $\mathrm{d}\log$ absorption of
Remark~\ref{rem:bar-pole-absorption}: the bar propagator extracts
residues against $\mathrm{d}\log(z-w) = \mathrm{d}z/(z-w)$, so a
pole of order~$n$ in the OPE produces a residue of order $n-1$ in
the collision expansion.

Part~(ii) is verified by the explicit computations for each standard
family, and the witnesses already force independence:
\[
\begin{array}{c|c|c|c}
\text{family} & p_{\max} & k_{\max} & r_{\max}\\
\hline
\cH_k & 2 & 1 & 2\\
V_k(\mathfrak{sl}_2) & 2 & 1 & 3\\
\beta\gamma_\lambda & 1 & 0 & 4\\
\mathrm{Vir}_c & 4 & 3 & \infty\\
\cW_N & 2N & 2N-1 & \infty .
\end{array}
\]
The first two rows have the same binary pole data and different
shadow depth. The last two rows have infinite shadow depth and
different generator pole order. The $\beta\gamma$ row is the decisive
contact witness: its generator OPE is simple-pole only, so
$p_{\max}(\beta\gamma)=1$ and
$k_{\max}(\beta\gamma)=0$; the composite
$J={:}\beta\gamma{:}$ has a self-OPE with a double contraction,
which contributes the quartic contact class recorded in the
landscape census and in
Remark~\ref{rem:class-C-beta-gamma-termination}. Hence
$r_{\max}(\beta\gamma)=4$.

Part~(iii) is the shadow-depth classification. Class~$G$ is
characterized by termination at degree~$2$ (all higher shadows
vanish); class~$L$ by termination at degree~$3$ (the Jacobi
obstruction is the last nonvanishing term); class~$C$ by termination
at degree~$4$ (the quartic contact class is the last nonvanishing
term); class~$M$ by non-termination (all degrees have nontrivial
shadow contributions). The identifications for each standard family
are the content of the landscape census
\textup{(}Chapter~\textup{\ref{ch:landscape-census}}\textup{)}.
\end{proof}

\section{The load-bearing role of the distinction}
\label{sec:three-invariants-load-bearing}

The distinction among $p_{\max}$, $k_{\max}$, and $r_{\max}$ is
load-bearing for three structural statements of this manuscript:

\begin{enumerate}
\item \emph{The operator-order trichotomy}
 \textup{(}Theorem~\textup{\ref{thm:shadow-depth-operator-order}}\textup{)}:
 the GZ\textup{26} commuting-Hamiltonian differential operator order
 is controlled by $k_{\max}$, not by $r_{\max}$. The three cases
 $k_{\max} = 0, 1, \geq 3$ produce genuinely different
 operator structures (trivial / multiplication / differential).

\item \emph{The shadow obstruction tower classification}
 \textup{(}the four classes $G, L, C, M$\textup{)}: the class is
 determined by $r_{\max}$, not by $p_{\max}$. In particular, class~$C$
 ($\beta\gamma$) has $r_{\max} = 4$ even though $p_{\max} = 1$.

\item \emph{The DNP non-renormalization theorem}
 \textup{(}Theorem~\textup{\ref{thm:non-renormalization-tree}}(iii)\textup{)}:
 one-loop exactness of line-operator OPE corresponds to $E_2$-collapse
 of the bar spectral sequence, which is a degree-$2$ statement
 controlled by the interaction of $p_{\max}$ and the Koszul property
 of~$\cA$, independently of $r_{\max}$.
\end{enumerate}

\begin{warning}[Never conflate the three finite invariants]
\label{warn:ap59-three-invariants}
A recurring error is to
conflate ``shadow depth'' (which means $r_{\max}$, the degree of the
obstruction tower) with ``collision depth'' (which means $k_{\max}$,
the pole order of the degree-$2$ collision residue) or with
``maximal OPE pole order'' (which means $p_{\max}$, the pole order
in the generator self-OPE). These are three distinct invariants
obeying the relation $k_{\max} = p_{\max} - 1$ but otherwise
independent. The $\beta\gamma$ system is the archetypal witness
of the independence: $(p_{\max}, k_{\max}, r_{\max}) = (1, 0, 4)$,
all three values distinct. Any statement of the form ``shadow depth
is the OPE pole order'' is \emph{wrong} in general.

Operational rule, applied at every cite: when a source asserts
``the depth of $\cA$ is $d$'', first ask
\emph{which} of the three invariants is meant; replace the
ambiguous symbol with one of $p_{\max}, k_{\max}, r_{\max}$
(or with $d_{\mathrm{gen}}, d_{\mathrm{alg}}$ for the depth-gap
invariants);
only then proceed. The most common
mismatches are: $p_{\max}$ vs.\ $k_{\max}$ (off by one through
$d\log$ absorption) and $k_{\max}$ vs.\ $r_{\max}$ (degree-$2$
vs.\ all-degree). The $\beta\gamma$ witness separates all
three.
\end{warning}

\begin{proposition}[Typed separation from scalar, centre, and bar--cobar data]
\label{prop:typed-separation-three-invariants}
\ClaimStatusProvedHere
Let $\cA$ be a completed standard-landscape chiral algebra on the
generic PBW/Koszul and perfect finite-window lanes on which
Theorems~\textup{C},~\textup{D}, and~\textup{H} are stated. The data
\[
(p_{\max},k_{\max},r_{\max}),\qquad
\kappa(\cA),\qquad
Z^{\mathrm{der}}_{\mathrm{ch}}(\cA)=\ChirHoch^\bullet(\cA,\cA),
\qquad
\bar B^{\mathrm{ch}}(\cA)\to \cA^{\mathrm i}\xrightarrow{\mathbb{D}_X}\cA^!
\]
belong to four different type signatures. The only universal
implications used in this chapter are:
\begin{enumerate}[label=\textup{(\roman*)}]
\item $p_{\max}$ determines $k_{\max}$ by
 $k_{\max}=p_{\max}-1$ on the chosen minimal strong generators;
\item $r_{\max}$ determines the coarse shadow class
 $G/L/C/M$ on the completed non-critical rows;
\item $\kappa(\cA)$ enters Theorem~D only through the scalar
 projection $\mathrm{obs}_g(\cA)=\kappa(\cA)\lambda_g$ on the
 uniform-weight scalar-diagonal lane, and through the genus-$1$
 all-weight scalar;
\item Theorem~C records the scalar Verdier sum
 $K^\kappa(\cA)=\kappa(\cA)+\kappa(\cA^!)$ after the Verdier partner
 and anomaly-ratio hypotheses have been fixed; it is not an
 equivalence of naive centres, derived centres, or brace
 dg~algebras;
\item Theorem~H is the cohomological-amplitude statement for the
 chiral Hochschild complex under the PBW/perfectness hypotheses; it
 is not a statement about the bare bar coalgebra.
\end{enumerate}
Consequently no scalar equality, no value of $r_{\max}$, and no
triple $(p_{\max},k_{\max},r_{\max})$ recovers the full derived centre
or the full Koszul bar--cobar package.
\end{proposition}

\begin{proof}
The type signatures are distinct by construction. The bar complex is
the conilpotent coalgebra $T^c(s^{-1}\bar\cA)$ with its chiral
differential; its cohomology $\cA^{\mathrm i}$ is the Koszul-dual
coalgebra, and $\cA^!$ is obtained only after Verdier duality on
finite-type windows. The derived chiral centre is instead the chiral
Hochschild complex
$\ChirHoch^\bullet(\cA,\cA)$, whose Theorem~H concentration uses the
ordered Fulton--MacPherson collapse followed by the PBW averaging
hypothesis. The scalar $\kappa$ is the genus-$1$ curvature trace of
the universal obstruction; Theorem~D promotes it to
$\kappa\lambda_g$ only on the uniform-weight scalar-diagonal lane,
while multi-weight all-genus free energies include
$\delta F_g^{\mathrm{cross}}$.

The low-degree witnesses exclude the possible collapses. Heisenberg
and affine Kac--Moody have the same
$(p_{\max},k_{\max})=(2,1)$ but different shadow depths
$2$ and~$3$. The $\beta\gamma$ system has
$(p_{\max},k_{\max},r_{\max})=(1,0,4)$, scalar complementarity
$K^\kappa(\beta\gamma)=0$ with its $bc$ partner, and nonzero
standard-family quartic shadow $S_4=-5/12$; thus scalar cancellation
does not annihilate the shadow tower and does not identify derived
centres. For $\mathcal W_3$, the diagonal scalar term is
$\kappa(\mathcal W_3)\lambda_g^{\mathrm{FP}}$, but at genus~$2$ the
full all-weight free energy has the nonzero cross-channel correction
recorded in Theorem~\ref{thm:multi-weight-genus-expansion}. These
examples prove that the arrows listed in the statement are the
available arrows, and no inverse arrow is available without the extra
finite-window, PBW, Verdier, or perfectness hypotheses named there.
\end{proof}

\section{The trichotomy at degree~2}
\label{sec:three-invariants-trichotomy}

On the bosonic integer-weight standard rows used by the GZ\textup{26}
operator-order comparison, the classification by $k_{\max}$ produces a
\emph{trichotomy}, not a dichotomy, reflecting three structurally
different behaviours of the GZ\textup{26} commuting Hamiltonians:

\begin{center}
\renewcommand{\arraystretch}{1.3}
\begin{tabular}{ccll}
$k_{\max}$ & \textbf{Regime} & \textbf{GZ\textup{26} $H_i$ are} & \textbf{Examples} \\
\hline
$0$ & trivial & trivial (no commuting structure at degree 2) & $\beta\gamma$ \\
$1$ & multiplication & scalar / matrix multiplication operators
 & Heisenberg, affine KM \\
$\geq 3$ & differential & genuine differential operators of order $k_{\max}-1$
 & Virasoro, $\Walg_N$ \\
\end{tabular}
\end{center}

\begin{remark}[Why $k_{\max} = 2$ is missing]
\label{rem:k-max-2-missing}
The value $k_{\max}=2$ is absent from the selected standard rows of
the operator-order theorem, not from the universe of all integer-weight
vertex algebras. It would require a genuine cubic generator-pair pole
$p_{\max}=3$. None of the rows used there has such a maximal binary
pole: $\beta\gamma$ has a simple mixed generator pole, Heisenberg and
affine Kac--Moody have quadratic current poles, and the Virasoro /
principal $\cW_N$ rows have maximal self-OPE poles at least quartic.
A chiral algebra with a cubic maximal generator-pair pole would define
an additional comparison row; it is not part of
Theorem~\ref{thm:shadow-depth-operator-order}.
\end{remark}

\begin{theorem}[The $k_{\max}$ trichotomy; \ClaimStatusConditional]
\label{thm:k-max-trichotomy}
For every chiral algebra $\cA$ in the bosonic integer-weight standard
sublandscape used by Theorem~\textup{\ref{thm:shadow-depth-operator-order}},
$k_{\max}(\cA)\in\{0,1\}\cup\{3,4,5,\ldots\}$. The value $2$ is
absent from that comparison table \textup{(}Remark
\textup{\ref{rem:k-max-2-missing})}. The classification by
$k_{\max}$ matches the GZ\textup{26} Hamiltonian operator-order
trichotomy of
Theorem~\textup{\ref{thm:shadow-depth-operator-order}}.
\end{theorem}

\begin{proof}
The bound follows by reading the minimal-generator OPE tables of the
selected rows: $\beta\gamma$ gives $p_{\max}=1$, Heisenberg and
affine Kac--Moody give $p_{\max}=2$, and Virasoro/principal
$\cW_N$ give $p_{\max}\ge4$. The identity
$k_{\max}=p_{\max}-1$ converts these values to
$k_{\max}\in\{0,1\}\cup\{3,4,5,\ldots\}$.
The match with the operator-order trichotomy is
Theorem~\ref{thm:shadow-depth-operator-order}.
\end{proof}

\section{The canonical fingerprint}
\label{sec:canonical-fingerprint}

The three invariants $p_{\max}$, $k_{\max}$, $r_{\max}$ separate
$\beta\gamma$ from Heisenberg, and Virasoro from $\Walg_3$, but they
do not separate $\Walg_3$ from super-$\Walg_3$, nor the symplectic
boson from the symplectic fermion: two algebras with identical
$(p_{\max}, k_{\max}, r_{\max})$ can have inequivalent bar complexes.
The refinement is two additional slots: the Hodge-filtered
supertrace~$\chi_{\mathrm{VOA}}$, which records parity; and the
\emph{coset slot}, which records the isomorphism class of
$\bigl(\mathrm{Com}(\cH_\kappa, \cA),\, \cA/\mathrm{Com}\bigr)$
and, at~$\kappa = 0$, the Feigin--Frenkel coset~$\mathrm{coset}_{FF}$.
Together with a count of minimal strong generators, these five slots
form the \emph{canonical fingerprint}. The fingerprint is a typed
projection from the modular-Koszul package; it is not the package
itself. The content of this section is the triangle
\begin{equation}\label{eq:fingerprint-triangle}
\text{modular-Koszul package}\;\overset{\text{(I)}}{\longrightarrow}\;\text{fingerprint}\;\overset{\text{(II)}}{\longrightarrow}\;\text{shadow class},
\end{equation}
where~(I) is the finite list of typed projections below and~(II) is
a surjective coarse projection (five-class stratum). The quadrichotomy
$G, L, C, M$
of~\S\ref{sec:three-invariants-relations} is the image of~(II)
restricted to~$\kappa \neq 0$; the fifth class~$\mathrm{FF}$ is the
locus~$\kappa = 0$.

\begin{definition}[Canonical fingerprint]
\label{def:canonical-fingerprint}
\index{canonical fingerprint $\varphi(\cA)$|textbf}
\index{fingerprint@$\varphi(\cA)$|see {canonical fingerprint}}
Let~$\cA$ be a chiral algebra with a chosen minimal strong generating
set. The \emph{canonical fingerprint} of~$\cA$ is the $5$-tuple
\begin{equation}\label{eq:fingerprint-def}
\varphi(\cA) \;:=\; \bigl(p_{\max}(\cA),\; r_{\max}(\cA),\; \chi_{\mathrm{VOA}}(\cA),\; n_{\mathrm{strong}}(\cA),\; \mathrm{coset}(\cA)\bigr) \;\in\; \mathcal{F},
\end{equation}
where the five slots are:
\begin{enumerate}[label=\textup{(\arabic*)},leftmargin=*]
\item $p_{\max}(\cA)$, the generator OPE pole order
\textup{(}Definition~\textup{\ref{def:p-max}}\textup{)};
\item $r_{\max}(\cA) \in \{2, 3, 4, \infty\}$, the shadow depth
\textup{(}Definition~\textup{\ref{def:r-max}}\textup{)};
\item $\chi_{\mathrm{VOA}}(\cA) := \operatorname{str}_{F^0\!\cA}(q^{L_0})$,
the Hodge-filtered supertrace, a formal power series in~$q$ with
$\mathbb{Z}$ coefficients;
\item $n_{\mathrm{strong}}(\cA) \in \mathbb{Z}_{\geq 1} \cup \{\infty\}$,
the cardinality of the chosen minimal strong generating set;
\item $\mathrm{coset}(\cA)$, the isomorphism class of the pair
\[
\bigl(\mathrm{Com}(\cH_\kappa, \cA),\;\cA / \mathrm{Com}(\cH_\kappa, \cA)\bigr) \quad\text{when } \kappa(\cA) \neq 0,
\]
and $\mathrm{coset}(\cA) := \mathrm{coset}_{FF}(\cA)$, the
Feigin--Frenkel coset class at the critical level, when~$\kappa(\cA) = 0$.
\end{enumerate}
Slot~(4) is invariant under automorphism of the generating set
\textup{(}Remark~\textup{\ref{rem:p-max-generator-choice}}\textup{)};
the collision depth $k_{\max}$ is determined by slot~(1)
via~\eqref{eq:k-max-from-p-max} and is therefore absent from~$\varphi$.
\end{definition}

\begin{theorem}[Fingerprint projection]
\label{thm:fingerprint-completeness}
\ClaimStatusProvedHere
For~$\cA$ in the completed standard landscape with a chosen minimal
strong generating set, the fingerprint factors through the typed
projection
\begin{equation}\label{eq:fingerprint-completeness}
\Pi_{\mathrm{fp}}\colon
\bigl(\bar B^{\mathrm{ch}}(\cA),\Theta_\cA,
Z^{\mathrm{der}}_{\mathrm{ch}}(\cA),\kappa(\cA),
\mathrm{coset}(\cA)\bigr)
\longrightarrow
\varphi(\cA),
\end{equation}
whose five components are:
\[
\pi_{\mathrm{pole}},\quad
\pi_{\mathrm{sh}},\quad
\pi_{\mathrm{str}},\quad
\pi_{\mathrm{gen}},\quad
\pi_{\mathrm{coset}}.
\]
They read respectively the maximal generator pole order, the maximal
nonzero shadow projection, the Hodge-filtered supertrace, the number
of minimal strong generators, and the commutant/critical coset slot.
If two typed modular-Koszul packages are isomorphic in the finite
PBW windows and their shadow projections and coset slots agree, then
their fingerprints agree. The converse is not asserted: equality of
fingerprints determines the recorded finite-row data and the coarse
shadow class, but not the full OPE, not
$Z^{\mathrm{der}}_{\mathrm{ch}}$, and not
$\bar B^{\mathrm{ch}}$ as an $\mathrm{Ass}^{\mathrm{ch}}$-coalgebra.
\end{theorem}

\begin{proof}
The five maps are read in different degrees and from different
objects:
\[
\begin{tikzcd}[row sep=small, column sep=large]
\bar B^{\mathrm{ch},2}(\cA) \arrow[r, "\pi_{\mathrm{pole}}"] &
p_{\max}(\cA) \arrow[r] & k_{\max}(\cA) \\
\Theta_\cA \arrow[r, "\pi_{\mathrm{sh}}"] &
r_{\max}(\cA) \arrow[r] & G/L/C/M \\
\cA \arrow[r, "{(\pi_{\mathrm{str}},\pi_{\mathrm{gen}})}"] &
(\chi_{\mathrm{VOA}},n_{\mathrm{strong}}) \\
\mathrm{Com}(\cH_\kappa,\cA) \arrow[r, "\pi_{\mathrm{coset}}"] &
\mathrm{coset}(\cA).
\end{tikzcd}
\]
The first line uses only the binary bar differential and the
$d\log$ absorption formula. The second line uses the shadow
obstruction tower, not the scalar trace of that tower. The third line
uses the graded vector space with Hodge filtration and the chosen
minimal strong generators. The fourth line uses the commutant or, at
critical level, the Feigin--Frenkel coset. An isomorphism of typed
packages preserves each source object and therefore preserves the five
read-offs.

The non-converse is forced by
Proposition~\ref{prop:typed-separation-three-invariants}. The scalar
Verdier sum of Theorem~C, the scalar genus trace of Theorem~D, the
chiral Hochschild complex of Theorem~H, and the bar--cobar package
have different domains and codomains. Forgetting to the five slots of
$\varphi$ cannot reconstruct the maps, higher braces, ordered
coinvariants, or all OPE coefficients that were forgotten. On a fixed
completed row of the landscape census with its normal form already
chosen, the five slots recover exactly the finite list of row data
recorded there;
that row-level recovery is the only reconstruction used in the
coarse projection below.
\end{proof}

\begin{remark}[Functorial symmetry under Koszul duality]
\label{rem:fingerprint-koszul-symmetry}
On rows whose Verdier partner is fixed in the landscape census, the
fingerprint projection is Koszul-equivariant:
$\varphi(\cA^!)$ is obtained from~$\varphi(\cA)$ by
$(p_{\max}, r_{\max})$ preserved,
$\chi_{\mathrm{VOA}} \mapsto \chi_{\mathrm{VOA}}^\vee$
(the Verdier-dual character),
$n_{\mathrm{strong}}$ preserved, and the coset slot exchanged under
the Koszul involution on the pair~$(\mathrm{Com}, \cA/\mathrm{Com})$.
This is functoriality of the projection
$\Pi_{\mathrm{fp}}$, not a reconstruction of the full bar coalgebra
or the derived centre from~$\varphi$.
\end{remark}

\begin{theorem}[Five-class stratum]
\label{thm:five-class-stratum}
\ClaimStatusProvedHere
The map
\[
 \varphi \longmapsto
 (r_{\max},\,\mathrm{coset},\,\kappa\text{-vanishing})
\]
stratifies the fingerprint space~$\mathcal{F}$ into five strata. The
Feigin--Frenkel stratum is the critical-level coset/centre stratum,
not the set-theoretic vanishing locus of the scalar~$\kappa$ alone:
\begin{center}
\renewcommand{\arraystretch}{1.25}
\begin{tabular}{lccl}
\textbf{Class} & $r_{\max}$ & $\kappa$ & \textbf{Archetype} \\
\hline
$G$ & $2$ & $\neq 0$ & Heisenberg $\cH_k$, lattice $V_\Lambda$ \\
$L$ & $3$ & $\neq 0$ & affine $\widehat{\fg}_k$ at~$k \neq -h^\vee$ \\
$C$ & $4$ & $\neq 0$ & $\beta\gamma$, $bc$ \\
$M$ & $\infty$ & $\neq 0$ & $\mathrm{Vir}_c$, $\Walg_N$ at $c \neq 0$ \\
$\mathrm{FF}$ & inherited & $0$ plus $\mathrm{coset}_{FF}$ & Feigin--Frenkel locus: $V_{-h^\vee}(\fg)$, critical $\Walg_N$
\end{tabular}
\end{center}
Every chiral algebra in the standard landscape lies in exactly one stratum.
\end{theorem}

\begin{proof}
The partition by~$r_{\max}\in\{2,3,4,\infty\}$ is
Proposition~\ref{prop:three-invariants-relations}(iii). The
critical-level refinement separates the Feigin--Frenkel locus from
the non-critical rows: at the affine critical level
$k=-h^\vee$, the Sugawara topologisation fails, the conformal vector
degenerates, the coset slot switches from
$(\mathrm{Com}(\cH_\kappa,\cA),\cA/\mathrm{Com})$ to
$\mathrm{coset}_{FF}$, and the chiral algebra acquires the
Feigin--Frenkel centre~$Z(V_{-h^\vee}(\fg))$. The implication
\[
 \mathrm{coset}(\cA)=\mathrm{coset}_{FF}
 \Longrightarrow \kappa(\cA)=0
\]
holds on the affine and critical $\cW$ lanes. The converse is not
used: if a different family has an accidental scalar value
$\kappa(\cA)=0$, its stratum is still determined by
$(r_{\max},\mathrm{coset})$. The five classes are mutually exclusive
because the pair~$(r_{\max},\mathrm{coset})$ takes five distinct
values on the landscape, and exhaustive by the landscape census
(Chapter~\ref{ch:landscape-census}).
\end{proof}

\begin{proposition}[Coarse projection functor]
\label{prop:coarse-projection-functor}
\ClaimStatusProvedHere
The projection
\begin{equation}\label{eq:coarse-projection}
\Pi_{\mathrm{coarse}} \colon \mathcal{F} \twoheadrightarrow \{G, L, C, M, \mathrm{FF}\},
\qquad
\Pi_{\mathrm{coarse}}\bigl(\varphi(\cA)\bigr) \;=\; \mathrm{class}(\cA),
\end{equation}
is functorial with respect to shadow-preserving morphisms of chiral
algebras \textup{(}Koszul duality, direct sum, central extension at
fixed level\textup{)}. On the open stratum~$\kappa \neq 0$, the
restriction~$\Pi_{\mathrm{coarse}}|_{\kappa \neq 0}$ lands in
$\{G, L, C, M\}$.
\end{proposition}

\begin{proof}
Functoriality on each stratum: direct sum adds shadow classes by
join (Proposition~\ref{prop:shadow-semilattice-standalone}), Koszul
duality preserves the class (Remark~\ref{rem:fingerprint-koszul-symmetry})
except on the~FF locus where~$\cA \leftrightarrow \cA^!$ exchanges
the two critical-level cosets. The restriction to~$\kappa\neq 0$
excludes the~FF stratum by definition.
\end{proof}

\begin{corollary}[Quadrichotomy $G/L/C/M$ as lossy projection]
\label{cor:quadrichotomy-as-projection}
The $G/L/C/M$ quadrichotomy of
Proposition~\textup{\ref{prop:three-invariants-relations}}(iii)
is the composition
\begin{equation}\label{eq:quadrichotomy-composition}
\mathrm{quat} \;=\; \Pi_{\mathrm{coarse}} \circ \varphi
\quad\text{restricted to}\quad \{\cA : \kappa(\cA) \neq 0\},
\end{equation}
obtained by discarding slots~\textup{(1), (3), (4), (5)} of the
fingerprint and retaining only the shadow depth~$r_{\max}$.
\end{corollary}

\begin{proof}
Restriction to~$\kappa \neq 0$ removes the~FF class; retention
of~$r_{\max}$ alone selects the stratum among the remaining four;
comparison with Proposition~\ref{prop:three-invariants-relations}(iii)
is term-by-term.
\end{proof}

\subsection{Witnesses of independence}
\label{ssec:fingerprint-witnesses}

The independence of the five slots is witnessed by standard pairs and
families in which one slot changes while the stated coarse slots are
held fixed. The witnesses are:

\begin{example}[Slot 1 vs.\ slot 2: $(p_{\max}, r_{\max})$ independence]
\label{ex:fingerprint-p-r-independence}
The pair $(p_{\max}, r_{\max})$ achieves
$(2, 2)$ for Heisenberg~$\cH_k$;
$(2, 3)$ for affine~$\widehat{\fg}_k$;
$(1, 4)$ for~$\beta\gamma$;
$(4, \infty)$ for~$\mathrm{Vir}_c$;
$(6, \infty)$ for~$\Walg_3$;
and $(2N, \infty)$ for~$\Walg_N$. The triples $(2, 2)$ and~$(2, 3)$
show that~$p_{\max} = 2$ does not determine~$r_{\max}\in\{2,3\}$;
the triples $(1, 4)$ and~$(4, \infty)$ show that the inequality
$p_{\max} \lessgtr r_{\max}$ can go in either direction.
\end{example}

\begin{example}[Slot 3: super-$\Walg_3$ vs.\ $\Walg_3$]
\label{ex:super-w3-vs-w3}
The algebras~$\Walg_3$ and super-$\Walg_3$ share
$(p_{\max}, r_{\max}) = (6, \infty)$ and the same number of strong
generators. They are separated by~$\chi_{\mathrm{VOA}}$ and the coset
slot:
\[
\chi_{\mathrm{VOA}}(\Walg_3) \;\in\; \mathbb{Z}[\![q]\!]_{+}, \qquad
\chi_{\mathrm{VOA}}(\text{super-}\Walg_3) \;\in\; \mathbb{Z}[\![q]\!]_{+} - q^{3/2}\mathbb{Z}[\![q]\!]_{+},
\]
with the odd fermionic generator of weight~$3/2$ entering the
supertrace with a minus sign. The coset~$\mathrm{Com}(\cH, \cdot)$
is bosonic in one case and super-bosonic in the other. Separating
these two algebras by fingerprint alone requires slot~(3).
\end{example}

\begin{example}[Symplectic boson vs.\ symplectic fermion: sign and coset]
\label{ex:symplectic-boson-vs-fermion}
The symplectic boson ($\beta\gamma$ at $\lambda = 1/2$, $c = -1$)
and the symplectic fermion~$\mathcal{SF}$ ($c = -2$) both belong to
class~$C$: $r_{\max} = 4$ with quartic contact invariant.
Slots~(1) and~(2) agree:
$(p_{\max}, r_{\max}) = (1, 4)$ in both cases. They are separated by
the sign of~$\chi_{\mathrm{VOA}}$ and by the coset:
\[
\chi_{\mathrm{VOA}}(\text{symp.\ boson})\;=\;+\mathrm{ch}(\beta\gamma_{1/2}), \quad
\chi_{\mathrm{VOA}}(\mathcal{SF})\;=\;-\mathrm{ch}(\mathcal{SF}),
\]
and the coset slot records~$\mathrm{Sp}(2)$ for the symplectic boson
(abelian symplectic automorphism group) versus
$\mathrm{OSp}(1|2)$ for the symplectic fermion (orthosymplectic,
involving the odd direction). Neither slot~(3) nor slot~(5) alone
suffices: their combination supplies the full separation.
\end{example}

\begin{remark}[Disambiguation: $k_{\max}$-trichotomy vs.\ $r_{\max}$-quadrichotomy]
\label{rem:trichotomy-vs-quadrichotomy}
Two distinct classifications live in this chapter and must not be
conflated. The $k_{\max}$-\emph{trichotomy}
\textup{(}Theorem~\textup{\ref{thm:k-max-trichotomy}}, three
operator-order regimes $0, 1, \geq 3$\textup{)} classifies
commuting-Hamiltonian behaviour at degree~$2$; its output has three
cells. The $r_{\max}$-\emph{quadrichotomy}
\textup{(}Proposition~\textup{\ref{prop:three-invariants-relations}}(iii),
four classes $G, L, C, M$\textup{)} classifies shadow-depth behaviour
across all degrees; its output has four cells. The five-class stratum
\textup{(}Theorem~\textup{\ref{thm:five-class-stratum}}\textup{)}
extends the quadrichotomy by the~FF locus. None of these three
classifications refines another on the nose: the trichotomy groups
Virasoro with~$\Walg_N$ ($k_{\max} \geq 3$), whereas the quadrichotomy
places them together in class~$M$ but separates~$\beta\gamma$
($k_{\max} = 0$, class~$C$) from Heisenberg ($k_{\max} = 1$, class~$G$).
The canonical fingerprint subsumes all three: the trichotomy is read
from slot~(1), the quadrichotomy from slot~(2), and the~FF locus from
slot~(5) via~$\kappa$-vanishing.
\end{remark}

\section{Compute-layer verification}
\label{sec:three-invariants-compute}

The three-invariants distinction is checked in the compute layer
by the engine
\texttt{compute/lib/theorem\_three\_paper\_intersection\_engine.py},
which contains the dictionary
\begin{verbatim}
 shadow_depth_to_operator_order('heisenberg')
 = {shadow_class: G, r_max: 2, p_max: 2, k_max: 1, op_order: 0}
 shadow_depth_to_operator_order('km_sl2')
 = {shadow_class: L, r_max: 3, p_max: 2, k_max: 1, op_order: 0}
 shadow_depth_to_operator_order('betagamma')
 = {shadow_class: C, r_max: 4, p_max: 1, k_max: 0, op_order: 0}
 shadow_depth_to_operator_order('virasoro')
 = {shadow_class: M, r_max: inf, p_max: 4, k_max: 3, op_order: 2}
 shadow_depth_to_operator_order('w3')
 = {shadow_class: M, r_max: inf, p_max: 6, k_max: 5, op_order: 4}
\end{verbatim}
The corresponding test file
\texttt{compute/tests/test\_theorem\_three\_paper\_intersection\_engine.py}
contains the critical
\texttt{test\_betagamma\_class\_C}
which verifies the $(p_{\max}, k_{\max}, r_{\max}) = (1, 0, 4)$
triple for the $\beta\gamma$ system.
